% Generated draft: requires mathematical and source review.
\documentclass[11pt]{article}
\usepackage[margin=25mm]{geometry}
\usepackage{fontspec}
\setmainfont{Latin Modern Roman}
\usepackage{amsmath,amssymb,mathtools,unicode-math}
\setmathfont{Latin Modern Math}
\usepackage{enumitem,needspace,xurl,hyperref}
\setlength{\parindent}{0pt}
\setlength{\parskip}{5pt}
\setlength{\emergencystretch}{3em}
\newcommand{\N}{\mathbb{N}}
\newcommand{\Z}{\mathbb{Z}}
\newcommand{\Q}{\mathbb{Q}}
\newcommand{\R}{\mathbb{R}}
\newcommand{\C}{\mathbb{C}}
\newcommand{\F}{\mathbb{F}}
\newcommand{\A}{\mathbb{A}}
\newcommand{\PP}{\mathbb{P}}
\newcommand{\E}{\mathbb{E}}
\newcommand{\eps}{\varepsilon}
\newcommand{\dd}{\,\mathrm{d}}
\begin{document}
\section*{UnsolvedMath ID: 30002373}
\textbf{Source problem code:} catalogue entry
\subsection{Definitions and mathematical statement}
Let \(\mathbb T\) denote the one-dimensional spatial torus, whose period is not specified in the supplied record, and write \(u_x=\partial_xu\), \(u_{xxx}=\partial_x^3u\), and \(u_t=\partial_tu\). Consider the two displayed families of equations with periodic spatial boundary conditions:
\[
 u_t+u_{xxx}-\partial_x(u^2)+N_4(x,u,u_x,u_{xx},u_{xxx})=0
\]
and
\[
 u_t+u_{xxx}\pm\partial_x(u^3)+N_4(x,u,u_x,u_{xx},u_{xxx})=0.
\]
The signs in the second display represent the two modified KdV equations. The underlying integrable equations are obtained by setting \(N_4=0\). The supplied hypothesis describes \(N_4\) as a Hamiltonian quasilinear perturbation vanishing to fourth order at the origin. Under a sufficiently smooth Taylor interpretation, with \(z=(u,u_x,u_{xx},u_{xxx})\), the order condition means
\[
 \partial_z^\alpha N_4(x,0)=0\qquad (|\alpha|<4).
\]
This interpretation does not require a nonzero fourth-order term. The record does not specify differentiability, uniformity in \(x\), the Hamiltonian structure, or the precise admissible quasilinear dependence; these remain unresolved hypotheses.

For clarity, conventional temporal recurrence notions can be described after choosing a normed space \(X\) of spatial functions. A trajectory \(t\mapsto u(t)\in X\), defined for all \(t\in\R\), is periodic if there exists \(T>0\) such that \(u(t+T)=u(t)\) for every \(t\). A finite-frequency quasiperiodic trajectory has a representation
\[
 u(t)=F(\theta+\omega t),\qquad F:(\R/\Z)^d\longrightarrow X,
\]
where \(d\geq1\), \(F\) is continuous, \(\theta\in(\R/\Z)^d\), and \(\omega\in\R^d\) satisfies \(m\cdot\omega\neq0\) for every \(m\in\Z^d\setminus\{0\}\). Under the Bohr interpretation, a continuous trajectory is almost periodic if, for every \(\varepsilon>0\), there exists \(\ell>0\) such that every interval of length \(\ell\) contains a \(\tau\) satisfying
\[
 \sup_{t\in\R}\|u(t+\tau)-u(t)\|_X<\varepsilon.
\]
These are explanatory conventions: the source does not specify the phase space, scalar field, solution regularity, or intended notion of almost periodicity, and it does not explicitly distinguish temporal recurrence from spatial periodicity.

\textbf{Persistence question.} For each displayed equation, do the periodic, quasiperiodic, or almost-periodic solutions of its integrable equation persist under sufficiently small Hamiltonian quasilinear perturbations \(N_4\) satisfying the stated order condition?

Here persistence is the target of the question, not a defined or established property. A fully quantified assertion would require specifying which unperturbed solutions are covered, the perturbation class and its smallness norm, the topology and limiting meaning of closeness, whether frequencies or periods may change, and the dependence of smallness thresholds on the solution. The record supplies none of these choices. In particular, it does not resolve whether every such solution is intended to persist or whether persistence of selected families is the intended target. Consequently, the displays and qualified question above preserve the supplied formulation without asserting an unsupported universal persistence theorem.
\subsection{Short English statement}
The problem asks whether recurrent solutions of integrable KdV and both signed modified KdV equations survive small Hamiltonian quasilinear perturbations whose nonlinear terms vanish to fourth order at the zero state. The precise solution class, perturbation size, recurrence convention, and meaning and quantifiers of persistence remain unspecified. The supplied dated curation labels the problem partially solved and reports KAM results for certain Cantor families; these are supplied assessments requiring independent verification.
\subsection{Sources}
\begin{itemize}
\item Local UnsolvedMath registry: \texttt{data/unsolvedmath/problems.json}, with \texttt{catalog/problems.json} as the supplementary registry; permanent UnsolvedMath ID \(30002373\), source problem code \texttt{catalogue entry}, title ``Persistence of KdV Solutions Under Quasilinear Perturbations.'' The supplied record belongs to the Oberwolfach Reports Open Problems collection. No individual proposer is identified; the proposed year is recorded as 2013.
\item Supplied report citation: \emph{Dynamische Systeme (2014)}, Oberwolfach Reports, pp.\ 1975--2033, DOI \texttt{10.4171/owr/2013/34}. \url{https://doi.org/10.4171/owr/2013/34}. The background describes these reports as records of mathematical workshops and their open-problem discussions. The citation's 2014 label and the 2013 DOI component and proposed year are retained without reconciliation.
\item Supplied statement-verification metadata labels the statement \texttt{corrected\_verified} and says that the two equations and the order condition were restored from the official report to replace a missing display reference. This is supplied provenance, not independent verification performed for this draft.
\item Baldi, Berti and Montalto, ``autonomous quasi-linear KdV'' (supplied reference label): \url{https://www.numdam.org/item/AIHPC_2016__33_6_1589_0/}.
\item Baldi, Berti and Montalto, ``forced KdV'' (supplied reference label): \url{https://arxiv.org/abs/1211.6672}.
\item Berti, Kappeler and Montalto, ``large KAM tori'' (supplied reference label): \url{https://arxiv.org/abs/1908.08768}.
\item Literature assessment checked at \texttt{2026-08-21}: supplied dated curation, not independently verified. The supplied status is \texttt{partially\_solved}. Its assessment states that KAM theory proves persistence and stability of Cantor families of small quasiperiodic solutions for broad Hamiltonian quasilinear KdV perturbations, while persistence of all periodic, quasiperiodic, or almost-periodic solutions under arbitrary small perturbations remains open. The background explicitly characterizes this as dated literature triage derived from supplied OWR curation and requests independent verification of time-sensitive status and cited claims. These claims are retained as attributed curation and do not establish their scope for both modified KdV signs.
\end{itemize}
\end{document}
